% Chapter 5: How to Describe Motion
\chapter{How to Describe Motion}

\step{A Counterintuitive Opening}

\begin{mainline}
A high-speed train sweeps across a plain at 300 kilometers per hour. Inside, a boy tosses an apple vertically upward. It traces an arc and falls back into his palm. It does not shoot toward the rear or hit the passenger behind him. Everything happens as though the train were standing still.

Something seems wrong. Think what a friend on the platform sees: during the half-second the apple is airborne, the train travels more than forty meters. Once released, the apple hangs in the air while the hand, seats, and floor rush along at over eighty meters per second. Why does it not fall behind? Who pushes it through the air to keep up with the train?

An equally strange thing happens on windows every day. Watch the side glass in rain: when the car is stopped, streaks run vertically downward; once it moves, all the streaks turn diagonal. The faster it goes, the more nearly horizontal they become. If rain falls downward, why does it move sideways on your window?

A larger oddity awaits. Earth carries us around the Sun at about 30 kilometers per second, over three hundred times a high-speed train's speed. The solar system in turn carries Earth around the Milky Way's center at about 220 kilometers per second. Every instant we travel faster than a cannonball could dream of, yet the coffee on the desk is motionless, a jump lands you where you started, and no morning of your life has brought a 220-kilometer-per-second wind.

Compare two experiences and the puzzle deepens: an airliner cruises at 900 kilometers per hour and you sleep soundly; a roller coaster barely exceeds a hundred and you scream yourself hoarse. Your body plainly has no interest in how fast you go. It is screaming at something else. Identifying that something is one of this chapter's central tasks.

These three oddities are really one. Behind them lies one of physics' simplest and deepest discoveries: \emph{motion is not a label attached to an object, but a relationship between an object and an observer}. This chapter will carefully redo the supposedly elementary task of describing motion until we can explain the apple, rain streaks, and Earth.
\end{mainline}

\step{Make a Guess First}

Before reading on, stake your intuition. Wrong guesses are more valuable than right ones, because intuition is what this chapter aims to repair.

Question one. Jump straight upward inside a train moving at constant velocity. Do you land (a) in the same place, (b) slightly behind it, or (c) slightly ahead? What if the train is \emph{braking} when you jump?

Question two. Two classmates argue. A says, ``I am sitting motionless on this high-speed train.'' B says, ``Nonsense, you are hurtling along at more than eighty meters per second.'' Who is right? Can you devise an experiment that settles it \emph{without looking out the window}?

Question three. Toss a ball straight upward. At the instant it reaches the top, what is its velocity? At that instant, is its tendency to accelerate, which we will call acceleration later, zero or nonzero?

Write your three answers down. We will check them at the chapter's end. If they resemble most people's answers, at least one has already put your intuition into a pit physicists took two thousand years to climb out of.

\step{Hands-on Experiment}

\begin{experiment}[Measuring a Fall with Slow-Motion Video]
\textbf{Materials}: a phone with slow-motion video, a tape measure or ruler, a coin or eraser, and a wall.

\textbf{Step one}: attach a paper strip vertically to the wall. Mark horizontal lines every 10 centimeters and label them 0, 10, 20, 30, and so on. This is your coordinate ruler.

\textbf{Step two}: release the coin from rest at the zero mark and film vertically in slow motion, preferably 240 frames per second, keeping the full scale in view.

\textbf{Step three}: play the video frame by frame. Call the release frame frame 0 and record the coin's marks at frames 5, 10, 15, and 20. In equal time intervals, the coin travels \emph{farther and farther}: its speed is not constant; it accelerates.

\textbf{Step four}: convert frame rate into time: at 240 frames per second, each frame is about \(1/240\) of a second. Convert frame numbers to elapsed seconds and plot position against time on graph paper. The points form a curve that grows steeper, not a straight line. Calculate how much the falling distance increases when falling time doubles. Most people measure close to fourfold: twice the time, four times the distance, precisely distance proportional to time squared.

\textbf{Record}: save the graph; later we will use it to measure a real constant of nature. If step three seemed to show constant-speed falling, you probably pressed pause: freely falling objects never move at constant speed.
\end{experiment}

Here is a free observation too: next time you ride a bus or subway moving at constant velocity, toss your keys straight up a short distance and catch them five times. Note any shift of the landing point relative to your hand. Then ask the driver, in imagination, to brake suddenly and repeat. Of course, do that version only in your head, or substitute a swing slowing down. Record both results for the ending.

\step{The Old Model Runs into Trouble}

First, give everyday language its due. We say ``This car is fast---a hundred and twenty kilometers per hour,'' never specifying relative to whom. That convention survives because we all share a default observation platform: the ground. Unless stated otherwise, speeds are ground-relative and everyone gets along. The old model is \emph{that position and velocity are an object's own properties; state them clearly and there is no need to mention an observer}.

Then difficulties arrive one after another.

Difficulty one: the illusion of neighboring trains. You sit in a stopped train while another on the adjacent track starts slowly. For a second or two you cannot tell whether it moves forward or you move backward. Looking down at the ground solves the mystery, but what if these were two parallel conveyor belts with no ground visible below? The old model calls who is moving an objective fact; your eyes say that watching only the other train cannot decide.

Difficulty two: who is speeding? Navigation warns that you have exceeded the limit by measuring your car relative to the \emph{ground}. To an overhead satellite, both you and the police car move at hundreds of meters per second with Earth's rotation; to the Sun, both circle at 30 kilometers per second. Fast and slow lose even their definitions without an observer.

Difficulty three: diagonal rain. A platform guard sees it fall vertically while you see it slant toward you from the train. The same rain, two trajectories, and \emph{both descriptions are right}. The old model asks which way rain really falls. Nature refuses to answer; it answers only relative to whom.

Difficulty four: does the Sun move? Every morning you see it rise above the horizon, a perfectly valid motion description in your frame. Copernicus says the same event can be described with the Sun stationary and the ground rotating downward under your feet, carrying you into sunlight. Both descriptions predict sunrise to the same second. The old model rushes to ask who is really right, but really is the redundant word: kinematics concerns relative relationships, not whose motion is authentic. Choosing a frame is about convenience, not truth. Of course, when dynamics asks why objects move this way, some frames will prove much more reasonable than others. That comes later.

Difficulty five: a philosophical hole. Someone proposes declaring a genuinely stationary judge's seat, with all velocities reported relative to it. Newton imagined exactly this and called it absolute space. The problem is Galileo's iron rule: in a sealed cabin moving uniformly, every mechanical experiment---throwing balls, dripping water, or swinging pendulums---has precisely the same result as in a stationary cabin. If absolute rest existed, such experiments should expose some trace of it, yet it never appears. A judge's seat that can never be measured is equivalent to nonexistent for physics.

The five difficulties share a diagnosis: the old model mistakenly treats motion as an object's own label. The escape is not a more authoritative judge, but admitting that \emph{position and velocity must include the observer from their very definitions}. This requires concepts obvious today but revolutionary when invented.
\step{Inventing New Concepts}

\begin{mainline}
\textbf{First concept: a reference frame.} To describe motion, specify three things: where an observer stands, the origin; the rulers used, the coordinate axes; and the clock used to time events. Together these form a \emph{reference frame}. Where an object is now has a strict meaning: a set of readings relative to a particular frame. What is position like? Latitude and longitude on a map: a place gets coordinates only after the prime meridian and equator are agreed. What is it unlike? A person's height, which does not change when you change maps, whereas position readings can be rewritten completely by changing frames.

\textbf{Second concept: displacement differs from distance traveled.} Distance is the total length actually covered; displacement is the directed segment from start to finish, Chapter 4's vector. Run one complete track lap and distance is 400 meters while displacement is zero. Taxi apps charge by distance because they care how much the tires rolled; missile guidance cares about displacement because it asks only how far and in which direction the target lies. Confusing them is kinematics' most elementary and common error.

You use position as a set of readings daily. An elevator's floor display is a one-dimensional coordinate system: origin at ground level, positive upward, and negative at the first basement. Coordinates can be negative, and nobody finds floor minus one strange. Navigation spreads a two-dimensional coordinate grid over a city and locates you at a point on it. Coordinate systems are not mathematicians' inventions; they simply clean up our existing habit of reporting locations.

\textbf{Third concept: velocity is the rate of position change, with direction.} Whoever changes position more in the same time moves faster. Crucially, distinguish \emph{average velocity}, averaged over the whole trip, from \emph{instantaneous velocity}, at a particular moment. The changing dashboard number is instantaneous velocity's magnitude. Bolt's 9.58-second hundred meters gives an average around 10.4 meters per second, but his later instantaneous speed approaches twelve, while it is zero at the start. One average hides the entire story.

Navigation's estimated arrival time is average-speed business: divide remaining distance by recent average speed, then adjust for traffic. Racing along at 120 kilometers per hour barely advances the estimate because the earlier ten-minute jam weighs down the average. Conversely, even stopping for coffee will not change the estimate if the overall average stays the same. Average and instantaneous velocities are each honest; the question is which one you use for which problem.

\textbf{Fourth concept: acceleration is the rate of velocity change.} Changing speed, by starting or braking, is acceleration; changing direction, by turning, is acceleration too. A bus starting makes you lean backward, braking makes you lean forward, and turning throws you sideways: three attacks by the same quantity. Acceleration is this chapter's deepest seed. Next chapter you will discover that force governs acceleration, while motion itself never needs maintaining.

\textbf{Fifth concept: state.} Combine where an object is now with how fast and in which direction it moves, and you obtain its \emph{state}. Why are both essential? Consider billiards: a cue ball at the same position produces entirely different next scenes when nudged gently or struck hard. Same position, different future; the difference is velocity. Given the state and a law, you can predict the next moment. This is the first appearance of our three recurring questions: what state is the system in, what law changes it, and how do we measure it?

\textbf{Sixth concept: the Galilean principle of relativity.} In all frames moving uniformly in straight lines without acceleration relative to one another, called inertial frames, mechanical experiments give identical results. The laws of falling and collision measured aboard a ship match those measured ashore. No experiment tells you whether the ship is really moving. In a sharper reverse formulation, \emph{uniform motion cannot be felt; only changes in motion can be detected}.
\end{mainline}

\begin{center}
\begin{tikzpicture}[scale=0.92]
  % Ground frame
  \draw[thick] (-5.4,0) -- (5.4,0);
  \fill[pattern=north east lines] (-5.4,-0.26) rectangle (5.4,0);
  \node[below=8pt] at (-4.6,0) {Ground frame};
  % Train
  \draw[thick,fill=gray!12] (-1.2,0) rectangle (3.6,1.3);
  \draw[thick,fill=white] (-0.9,0.35) circle (0.28);
  \draw[thick,fill=white] (2.9,0.35) circle (0.28);
  \node at (1.2,0.85) {Constant-velocity train};
  \draw[-{Stealth[length=3mm]},very thick,blue!70!black] (3.7,1.6) -- (5.3,1.6) node[right=1pt] {\(v\)};
  % Person and apple
  \draw[thick,fill=orange!20] (0.9,0.55) circle (0.16);
  \node[below=1pt] at (0.9,0.55) {\scriptsize Person};
  \draw[thick,fill=red!40] (1.9,0.75) circle (0.09);
  % From the train: vertical
  \draw[dashed,red!70!black] (1.9,0.55) -- (1.9,1.85);
  \node[red!70!black,above=0pt] at (1.9,1.85) {\scriptsize Train view: straight up and down};
  % From the ground: parabola
  \draw[dashed,blue!70!black] (1.9,0.75) parabola bend (2.7,1.7) (3.5,0.75);
  \node[blue!70!black,below=2pt] at (3.9,0.6) {\scriptsize Ground view: parabola};
\end{tikzpicture}
\end{center}

This diagram already previews the ending: the same apple moves straight up and down to someone aboard and follows a parabola to someone on the platform. Both descriptions are right and follow from one law. That is the power of a reference frame: it prevents the descriptions from fighting.

\step{The Model in One Sentence}

\begin{oneline}
Motion has no absolute label: state your reference frame before position and velocity have meaning; all uniformly moving observers are equals, and none can prove through a mechanical experiment that they are really moving.
\end{oneline}

Remember this and the train's apple, slanted rain, and racing Earth are already half explained. Mathematics will handle the other half.

\step{The Mathematical Translator}

\begin{mainline}
Translate velocity first. Write position as \(x\) and time as \(t\). A displacement \(\Delta x\) takes time \(\Delta t\), giving average velocity
\[
  \bar v=\frac{\Delta x}{\Delta t}.
\]
Read \(\Delta\) as change; the expression means position change divided by elapsed time. Routine checks: stationary means \(\Delta x=0\), hence \(\bar v=0\). Negative \(\Delta x\), motion the other way, makes \(\bar v\) negative and automatically preserves direction. Less time for the same displacement gives greater velocity, as intuition expects. Shrink \(\Delta t\) toward a single instant and average becomes instantaneous velocity. On the position--time graph from Chapter 3, average velocity is a secant's slope and instantaneous velocity a tangent's slope. \textbf{Slope means how fast change occurs}. This graph deserves another look:
\end{mainline}

\begin{center}
\begin{tikzpicture}[scale=1.05]
  \draw[-{Stealth[length=2.5mm]}] (0,0) -- (6.4,0) node[below left=2pt] {\(t\)};
  \draw[-{Stealth[length=2.5mm]}] (0,0) -- (0,3.6) node[below left=2pt] {\(x\)};
  % Constant-acceleration curve
  \draw[very thick,blue!70!black] (0,0) parabola (5.4,3.1);
  \node[blue!70!black,right=2pt] at (5.0,2.55) {Constant acceleration};
  % Constant-velocity line
  \draw[very thick,red!70!black] (0,0) -- (5.9,1.55);
  \node[red!70!black,right=2pt] at (5.9,1.55) {Constant velocity};
  % Schematic tangent
  \draw[dashed,gray] (2.4,0.47) -- (4.6,2.6);
  \node[gray,below=1pt] at (4.0,2.2) {\scriptsize Tangent slope at an instant \(=\) instantaneous velocity};
  \fill (3.5,1.19) circle (0.06);
\end{tikzpicture}
\end{center}

\begin{mainline}
Next translate acceleration:
\[
  a=\frac{\Delta v}{\Delta t},
\]
velocity change divided by elapsed time. Check: uniform motion gives \(\Delta v=0\), hence \(a=0\). However fast uniform motion is, it has no acceleration. On a train moving uniformly at 300 kilometers per hour, your acceleration is zero, which is why you cannot feel it. Reverse faster and faster, and velocity is negative and increasingly so: \(\Delta v<0\), with acceleration opposite the forward direction. The expression still works.

Uniform acceleration, constant \(a\), is the simplest and most useful case. With initial velocity \(v_0\), three formulas cover everything:
\[
  v=v_0+at,\qquad
  x=x_0+v_0 t+\tfrac12 a t^2,\qquad
  v^2=v_0^2+2a\,(x-x_0).
\]
Check each. Set \(a=0\): the first reduces to \(v=v_0\), constant velocity; the second to \(x=x_0+v_0 t\), uniform straight-line motion; the third to \(v=v_0\). They agree without acceleration. Set \(t=0\): the first gives \(v=v_0\), the second \(x=x_0\), both returning to the start. Reverse \(a\) for braking: the first predicts steadily decreasing velocity, then zero, then negative velocity. A stopped car does not actually reverse without reverse gear; this reminds us that the formula applies during braking, not afterward. All direction checks pass.

\textbf{An estimate using real data: braking distance.} On dry pavement, a car braking hard has acceleration magnitude about \(5\) meters per second squared, roughly half a \(g\). Its speed is \(50\) kilometers per hour, about \(13.9\) meters per second. Use the third formula with final velocity \(v=0\):
\[
  0=v_0^2-2a\,s\quad\Rightarrow\quad s=\frac{v_0^2}{2a}=\frac{13.9^2}{2\times5}\approx 19\ \text{m}.
\]
Double the speed to \(100\) kilometers per hour? Velocity appears \emph{squared}, so distance quadruples to about \(77\) meters, excluding another \(28\) meters during the driver's roughly one-second reaction. That is why twice the speed means roughly four times the danger, not twice. This formula is written on the reverse of highway signs telling you to keep your distance.

Use the same formulas backward to see how engineers \emph{deliberately} reduce acceleration. A high-speed train accelerating from rest to 300 kilometers per hour, about 83 meters per second, does not get full throttle instantly. The driver allows about five minutes for a leisurely increase, giving average acceleration
\[
  a=\frac{\Delta v}{\Delta t}=\frac{83\ \text{m/s}}{300\ \text{s}}\approx 0.28\ \text{m/s}^2,
\]
less than three percent of \(g\). This is why a cup of coffee can sit steadily inside. Passenger comfort effectively limits acceleration, not velocity. Roller coasters instead cultivate acceleration: almost all their thrill comes from sharp changes in its magnitude and direction, not speed itself. Theme parks make \(a\) large and railways small, both reading the same formula.

Your phone contains an acceleration surveyor. Its fingernail-sized accelerometer suspends a tiny mass that shifts slightly relative to the casing as the phone accelerates. Circuits translate the shift into three directional acceleration readings, enabling screen rotation, step counting, and tilt-controlled games. Curiously, resting on a table it reads not zero but an upward \(9.8\): it measures not change in motion itself, but how its support pushes it. This apparent contradiction will be fully resolved after forces in Chapters 6 and 7, and will return on a larger stage in Chapter 41.

Free fall is a special case of constant acceleration: near the ground, ignoring air resistance, all falling objects have the same downward acceleration \(g\approx 9.8\) meters per second squared. Return to your experiment: since the coin falls from rest, fitting \(x=\tfrac12 g t^2\) to your curve measures \(g\) personally. A one-meter fall, for example, predicts \(t=\sqrt{2\times1/9.8}\approx 0.45\) seconds. Check whether your slow-motion coin reaches the one-meter mark around frame 100 at 240 frames per second.
\end{mainline}

\begin{mathbridge}
Where do the three formulas come from? The mysterious \(\tfrac12\) in the second most deserves an explanation. Start at \(v_0\) and accelerate uniformly to \(v\). Since velocity grows \emph{uniformly}, its mean is exactly the mean of its endpoint values: \(\bar v=(v_0+v)/2\). Thus displacement is
\[
  \Delta x=\bar v\,t=\frac{v_0+(v_0+at)}{2}\,t=v_0 t+\tfrac12 a t^2.
\]
The \(\tfrac12\) was not inserted to make things work. It expresses the geometric fact that a uniformly increasing quantity's mean equals its endpoints' mean. On a velocity--time graph, displacement is the area between line and time axis: a rectangle plus a triangle, whose area naturally includes \(\tfrac12\). Eliminate \(t\) between the first two formulas to obtain the third. One level deeper, write the \(\Delta t\)-to-zero limit as derivatives: \(v=\dfrac{dx}{dt}\) and \(a=\dfrac{dv}{dt}\). Position, velocity, and acceleration form a chain of differentiation used next chapter.

Projectile motion can be settled here too. Resolve initial velocity horizontally and vertically using Chapter 4. With air resistance ignored, no horizontal push or pull acts, so \(v_{0x}\) stays constant. Vertically, only \(g\) acts, giving uniform acceleration. The directions proceed independently and simultaneously:
\[
  x=v_{0x}\,t,\qquad y=v_{0y}\,t-\tfrac12 g t^2.
\]
Eliminate \(t\) and \(y\) becomes a quadratic function of \(x\), a parabola. Check: a launch angle of \(90^\circ\), vertically upward, gives \(v_{0x}=0\) and leaves \(y=v_{0y}t-\tfrac12gt^2\), vertical constant acceleration. Let \(g\to0\), imagining gravity switched off, and \(y=v_{0y}t\) with \(x=v_{0x}t\) combines into a straight line. Without gravity, projectiles move straight, consistent with the first law. Velocity at the top is not zero: the vertical component is zero, but the horizontal component remains \(v_{0x}\).

This parabola is everywhere. After a basketball leaves your hand, horizontal motion is uniform while vertical speed first falls then rises; its center traces this curve. Fountain jets and fire-hose streams are parabolas drawn by successive water droplets. Why seek a higher release arc in basketball? A steeper descending parabola makes the hoop's effective opening larger from the ball's perspective, forgiving the same release-speed error more readily. Players may not know the formula, but their feel has calculated it thousands of times.
\end{mathbridge}

\begin{center}
\begin{tikzpicture}[scale=1.05]
  \draw[-{Stealth[length=2.5mm]}] (0,0) -- (7.0,0) node[below left=2pt] {\(x\)};
  \draw[-{Stealth[length=2.5mm]}] (0,0) -- (0,3.2) node[below left=2pt] {\(y\)};
  \draw[very thick,domain=0:6.4,smooth,variable=\t,blue!70!black] plot ({\t},{1.5*\t-0.195*\t*\t});
  % Initial velocity components
  \draw[-{Stealth[length=2.5mm]},very thick,red!70!black] (0,0) -- (1.7,1.28) node[above right=0pt] {\(v_0\)};
  \draw[-{Stealth[length=2.5mm]},thick,red!70!black] (0,0) -- (1.7,0) node[below=1pt] {\(v_{0x}\)};
  \draw[-{Stealth[length=2.5mm]},thick,red!70!black] (0,0) -- (0,1.28) node[left=1pt] {\(v_{0y}\)};
  \draw[dashed,gray] (1.7,0) -- (1.7,1.28) -- (0,1.28);
  % Highest point
  \fill (3.85,2.88) circle (0.07);
  \draw[-{Stealth[length=2.5mm]},thick,green!50!black] (3.85,2.88) -- (5.35,2.88) node[right=1pt] {\scriptsize Only horizontal velocity at the top};
  \draw[dashed,gray] (3.85,0) -- (3.85,2.88);
\end{tikzpicture}
\end{center}

\begin{mainline}
Finally translate the conversation between reference frames. A train moves at \(u\) relative to the ground, and you walk toward its front at \(v'\) relative to the train. The ground measures your velocity as
\[
  v=u+v'.
\]
This simple addition across frames is Galilean velocity addition. Check: walk toward the rear and \(v'\) is negative, giving \(v<u\). If you walk rearward at exactly \(v'=-u\), then \(v=0\): your friend on the ground sees you hovering in midair and retreating out of the station, while you feel you are merely strolling inside.

An escalator is this formula's cheapest laboratory. It moves upward at about half a meter per second; stand still on it and the ground measures half a meter per second. Walk upward at one meter per second relative to it and the ground measures one and a half. Turn around and walk downward at the same pace; if your pace exactly equals the escalator speed, you remain fixed in place. A nearby shop assistant sees you walking on the spot, though you feel you are working hard. The same legs, two frames, and two utterly different motion stories, neither a lie. Window rain follows the same formula: rain moves vertically relative to ground, the car horizontally relative to ground, and rain relative to the \emph{car} is their vector difference. The window coordinates therefore give vertical rain a horizontal component and sloping streaks, flatter at greater car speed. Opening oddity two is solved.
\end{mainline}

\begin{challenge}
Write changing frames as a full coordinate translation. Ground coordinates are \((x,t)\), and train coordinates move at \(u\) along \(x\). The Galilean transformation is
\[
  x'=x-ut,\qquad t'=t.
\]
Differentiate twice with respect to \(t\): \(v'=v-u\), then \(a'=a\). Notice the last result: \emph{velocity changes with reference frame, but acceleration does not change between inertial frames}. That is why we cannot casually decide what is relative and what everyone agrees on. Position and velocity are relative; acceleration and time are shared in Newton's world. Chapter 6's force is tied to acceleration, so Newton's laws look identical in every inertial frame. The mathematical heart of Galilean relativity lies in \(a'=a\).

Look ahead to a boundary. The transformation hides an unnoticed assumption: \(t'=t\), the two frames share one clock. It is faultlessly accurate in daily life, but experiments near light speed force us to modify it, the story of Chapter 38. Every formula here is a low-speed approximation, good enough for the whole of industrial civilization to rest upon.

Air resistance deserves an aside. Real falling motion roughly obeys \(a=g-kv/m\): greater velocity means more drag and less acceleration until velocity approaches \(v_{\text{terminal}}=mg/k\). Raindrops reach the ground at just such a constant terminal speed of a few meters per second. Your filmed coin remains close to pure free fall; replace it with a feather and this correction immediately becomes the main story. Models are brought closer to truth step by step in this way.
\end{challenge}

\step{The Model's Limits}

The conditions for this chapter's models deserve reading like operating instructions.

First, we treated objects as \emph{point particles}, points with mass but no shape. Describing a train's Beijing-to-Shanghai trip scarcely needs its two-hundred-meter length, so the model works well. Describing a gymnast's somersault immediately defeats it: different body parts move differently, requiring Chapter 11's rotational language. A model is suitable or unsuitable, rather than right or wrong.

Second, uniform-acceleration formulas require constant \(a\). Free fall near the ground approximates this, but roller-coaster acceleration changes continuously and \(v=v_0+at\) cannot be used directly even for one step. Divide the motion into tiny intervals, reversing the mathematical bridge's derivative chain, or use Chapter 9's energy methods.

Third, ignoring air resistance is an assumption, not a fact. Shot puts and coins comply well; shuttlecocks, raindrops, and parachutes do not. Galileo's claim that all objects fall equally fast applies in vacuum. Apollo 15 astronauts released a hammer and feather together on the Moon and they landed side by side, the law's most famous public performance.

Fourth, the ground is only \emph{approximately} an inertial frame. Earth rotates, so strictly we turn at every instant and acceleration is nonzero. Its effect is negligible in everyday life; how delicate an experiment reveals it is a calculation for later chapters.

Fifth, the easiest condition to forget: every speed here is assumed far below light speed. This is not an oversight but the identity card of all Newtonian mechanics.

Sixth, consider instantaneous itself. Every real measurement takes time: radar measures displacement over a brief interval, and a phone accelerometer averages a mass's deflection. Velocity at one instant is therefore never measured directly; it is the limit approached as the measurement window shrinks. Instantaneous values belong to models, averages to measurements. With a sufficiently small window they almost coincide, so daily use may mix them, but keep track of which is the idealization.

\step{Explaining the Opening Phenomenon}

Now close the cases.

\textbf{Why the apple does not fall behind.} Before the toss, it already travels with the train at more than eighty meters per second. Tossing adds only vertical velocity; nothing pushes or drags horizontally. By the principle of inertia, formally stated in Chapter 6, the horizontal velocity remains unchanged. Person, apple, and carriage move together horizontally, while vertically the apple makes an ordinary upward toss. Aboard, it goes straight up and down; from the platform, it follows a parabola advancing horizontally at over eighty meters per second. Galilean velocity addition translates the two descriptions exactly. The first opening answer is therefore the same place at constant velocity. While braking, the carriage slows but your horizontal velocity temporarily does not, so you land somewhat \emph{ahead}.

\textbf{Why rain moves sideways.} Streak direction is rain velocity relative to ground minus car velocity relative to ground. A faster car gives a larger horizontal component and flatter streaks.

\textbf{Why Earth's wind does not exist.} We, the atmosphere, and the coffee cup all move almost uniformly with Earth. Galilean relativity says uniform motion creates no locally detectable experimental effect; only \emph{change} can be felt. Earth's orbit is curved, so our motion is not strictly uniform, but its turning tendency is spread over one circuit a year and imperceptibly weak in daily life. The number 220 kilometers per second itself creates no sensation.

The complete opening answer follows: jumping vertically inside a uniformly moving train does not leave you behind because the jump changes only vertical motion. Horizontally you and the train were already together, with no mechanism to separate you. What really needs explaining is not why you stay with it, but why you thought you should fall behind: that is the old model's lingering illusion.

Settle the other two bets too. For question two, A and B are \emph{both right}: A reports in the carriage frame, B in the ground frame, and Galilean addition translates between them. Neither owns the real account. Without looking out the window and with access only to the cabin's internal happenings, no experiment can decide a winner: Galilean relativity's verdict. Question three remains a debt, hidden in challenge 1 below. If you can already answer instantly, the new concepts have taken root.

\step{Thought Experiments and Challenges}

\begin{thought}[Galileo's Windowless Cabin]
Imagine being below deck on a large ship with every door and window sealed and perfect insulation from sound and vibration. It moves uniformly in a straight line across a calm sea. You have balls, pendulums, water drops, spring scales, and balances. Can you design a mechanical experiment that measures the ship's speed?

Galileo's answer is no. Every cabin experiment---ball falling times, pendulum periods, drop shapes---matches a stationary ship. Imagine the cabin as Earth orbiting the Sun at thirty kilometers per second: we live in such a windowless cabin, and only when telescopes and the stars supplied a view outside did humanity become certain Earth was moving.

Push the scale further. A future spaceship cruises uniformly at 220 kilometers per second, like the solar system around the galaxy. Inside you drink coffee, toss balls, and weigh yourself exactly as in port. However large the speed, if it stays constant, no cell or instrument can know it. Nature seems to shout the same message: it concerns itself with change, not uniform motion.

One further question: if this principle is right, does light obey it? Will measuring light speed in the cabin give the same value as ashore? Three hundred years later, that innocent question shattered \(t'=t\) and forced special relativity into being, Chapter 38. Our model stands at that revolution's doorway.
\end{thought}

\begin{puzzles}
\item A vertically thrown ball has zero velocity at its highest point. Is its acceleration zero too? If acceleration truly were zero at that instant, what would happen next?\\
\emph{Hint}: acceleration concerns how fast velocity is changing, not its current value. Make a thought judgment: with zero acceleration at the top, what would this chapter's formulas predict for velocity next instant? The ball would remain suspended forever. Reality refuses.
\item Two cars travel on the same straight road, and their velocity--time curves intersect at some instant. Does that mean they meet, occupying the same position then?\\
\emph{Hint}: the graph records how fast, not where. An intersection means only equal velocities then. Position information lies in the \emph{area} between graph and time axis, with each starting position added.
\item A two-hundred-meter train moves uniformly at seventy meters per second. You walk from rear to front at two meters per second relative to the carriage. By the time you reach the front, how far has a platform observer measured you moving?\\
\emph{Hint}: first find the walking time, then your ground-relative velocity by Galilean addition. Alternatively, find how far the train's front moves during that time; you have traveled exactly one train length less than the front. Both routes must give the same number. If they disagree, something was calculated wrongly.
\item Two high-speed trains approach each other at 300 kilometers per hour each. As they pass, someone inside one sees the other flash by at a relative 600 kilometers per hour, approaching an airliner's speed. Why are both trains' passengers unharmed, with coffee barely trembling? Now consider one train hitting a wall coming toward it at 300 kilometers per hour: which is worse, this or hitting a stationary wall?\\
\emph{Hint}: harm comes from the severity of a forced change in motion, not a velocity reading relative to something passing by. Choose a passenger as the object of study and ask how much their velocity changes in an instant. As for the moving wall, who supplies its 300-kilometer-per-hour motion? If it can really drive toward you, how does it differ from a train?
\end{puzzles}
